\documentclass[a4paper,11pt]{article}
\usepackage[utf8]{inputenc}
\usepackage{lmodern}
\usepackage[T1]{fontenc}
\usepackage[brazil]{babel}
\usepackage[margin=1.5cm]{geometry}
\usepackage{hyperref}
\usepackage{enumitem}
\usepackage{titlesec}
\usepackage{xcolor}
\usepackage{parskip}

\hypersetup{
    colorlinks=true,
    linkcolor=blue,
    filecolor=magenta,      
    urlcolor=blue,
}

\titleformat{\section}{\large\bfseries\color{black}}{}{0em}{}[\titlerule]

\begin{document}

\begin{center}
    {\LARGE \textbf{Osmar Nunes Júnior}} \\ \vspace{0.2cm}
    Curitiba, Paraná // +55 (41) 9 9631-9647 \\
    \href{https://www.linkedin.com/in/osmarnunesjunior}{linkedin.com/in/osmarnunesjunior} \\
    \href{https://www.github.com/osmarjrn}{github.com/osmarjrn} \\
    \href{mailto:osmjr04@gmail.com}{osmjr04@gmail.com}
\end{center}

\vspace{0.5cm}

\section*{Da formação acadêmica}
Estudante de Tecnologia em Análise e Desenvolvimento de Sistemas na Universidade Federal do Paraná \textbf{(IRA: 0.812/1)}. Ensino Médio certificado pelo Colégio Militar de Curitiba \textbf{(MG: 8.5/10)}.

\section*{Das competências relevantes}
\textbf{Tecnologias}
\begin{itemize}[noitemsep, label={--}]
   \item \textbf{MySQL:} Banco de Dados I e II (2º e 3º período);
\item \textbf{Shell Script:} Administração de Sistemas (1º período);
\item \textbf{Java:} Linguagem de Programação Orientada a Objetos I (3º período);
\item \textbf{HTML, CSS, JavaScript e PHP:} Desenvolvimento Web I (2º período);
\item \textbf{Python:} Redes de Computadores (2º período) e Sistemas Operacionais (3º período);
\item \textbf{C, C++:} Linguagem de Programação (2º período) e Estrutura de Dados (3º período);
\item \textbf{Pacote Office (Excel, Word e PowerPoint):} Criação, edição e administração de documentos.
\end{itemize}

\vspace{0.3cm}
\textbf{Habilidades}
\begin{itemize}[noitemsep, label={--}]
  \item Informática avançada;
\item Edição de fotos e vídeos;
\item Suporte técnico a usuários;
\item Inteligência socioemocional;
\item Análise e resolução de problemas;
\item Manutenção e configuração de hardware;
\item Execução de estudos e adaptações técnicas;
\item Criação e gerenciamento de conteúdo digital;
\item Comunicação e negociação eficazes, e eficientes;
\item Redação de textos coesos e coerentes em norma culta;
\item \textbf {Idiomas:} Português nativo e Inglês intermediário (conversação e escrita).
\end{itemize}

\section*{Das experiências*}
\begin{itemize}[noitemsep, label={$\bullet$}]
  \item \textbf{Estagiário de TI no Centro RPG Curitiba}\\
\textit{Suporte e Desenvolvimento de TI (750 horas - 12/2025 até 06/2026)}\\
Protagonismo na análise, resolução e implantação de fichas e prontuários digitais.
\vspace{0.2cm}

\item \textbf{Membro da Empresa Júnior de Computação (ECOMP)}\\
\textit{Voluntário em Assessoria de Projetos (92 horas - 10/2025 até 04/2026)}\\
Administração de arquivos, com suporte à organização e gestão de projetos.
\vspace{0.2cm}

\item \textbf{Projeto Acadêmico: Desenvolvimento Full-Stack -- UFPR}\\
\textit{Tecnologia em Análise e Desenvolvimento de Sistemas (20 horas - 05/2026)}\\
Implementação em equipe de um jogo web interativo para digitação correta de cores.
\vspace{0.2cm}

\item \textbf{Suporte Técnico Independente}\\
\textit{Manutenção de Hardware e Software (Autônomo - paixão atemporal)}\\
Consultoria, configuração e manutenção de equipamentos com soluções tecnológicas.
\end{itemize}

\section*{Savoir-faire*}

\hspace*{1.25cm}Acadêmico de Tecnologia em Análise e Desenvolvimento de Sistemas na Universidade Federal do Paraná, com trajetória marcada pela disciplina atrelada a valores do exército brasileiro e pelo engajamento técnico desde a educação básica no Colégio Militar de Curitiba. Minha atuação inicial foi pautada pela participação ativa em clubes e atividades extracurriculares voltadas à Tecnologia da Informação e Relações Internacionais, onde desenvolvi as bases do pensamento analítico e colaborativo.

\hspace*{1.25cm}Assim, apliquei conhecimentos teóricos em ambiente corporativo no Centro RPG Curitiba, clínica referência de fisioterapia, onde liderei a transição de processos legados para o ecossistema digital. Entre as principais contribuições, destacam-se a implementação de prontuários digitais baseados em modelos físicos e o desenvolvimento de scripts de automação em Python para otimizar o fluxo de implantação de dados. Além da vertente de desenvolvimento, possuo experiência prática em suporte técnico, manutenção de infraestrutura de hardware e gerenciamento de sistemas internos.

\hspace*{1.25cm}No mais, dentro da Ecomp, atuei no suporte à gestão e estruturação de projetos, auxiliando na organização de documentações técnicas e no acompanhamento de cronogramas. Colaborei com a Diretoria de Projetos na facilitação de fluxos de trabalho e na garantia de qualidade das entregas, contribuindo para a organização interna e o alinhamento das equipes.

\hspace*{1.25cm}Enfim, tenho uma forte dedicação à análise e resolução de problemas, especialmente no contexto de tecnologia da informação, onde foco em encontrar soluções eficientes para desafios técnicos complexos. Como ponto de desenvolvimento, busco aprimorar constantemente minha gestão de tempo e o cumprimento de prazos, trabalhando para equilibrar demandas e prioridades.

\section*{Dos valores}
Disciplina, autonomia, iniciativa de atitude, liderança, espírito de corpo, resiliência e aprendizagem.

\section*{Do almejado}
Em busca de desenvolvimento e produção com oportunidade de contribuir efetivamente nesse exercício.

\end{document}
